\subsection{Cuerpos de definición de una parte de un preesquema.}

\begin{proposition}[4.9.1]
\label{IV.2.4.9.1-fr}
Sean $k$ un cuerpo, $X$ un $k$-preesquema, $K$ una extensión de $k$,
$\overline T$ una parte cerrada de $X_{(K)}$ y $K'$ una subextensión de $K$.
Las condiciones siguientes son equivalentes:
\begin{enumerate}
\item[\textup{a)}] $\overline T$ está definida sobre $K'$.
\item[\textup{b)}] La parte abierta $\overline U=X_{(K)}\setmin\overline T$ de $X_{(K)}$
está definida sobre $K'$.
\item[\textup{b')}] El subpreesquema de $X_{(K)}$ inducido sobre el abierto
$\overline U$ está definido sobre $K'$.
\item[\textup{c)}] Existe un subpreesquema cerrado $\overline Z$ de $X_{(K)}$ cuyo
espacio subyacente es $\overline T$ y que está definido sobre $K'$.
\end{enumerate}
\end{proposition}

\begin{proof}
Si $g:X_{(K)}\to X_{(K')}$ es la proyección canónica, decir que
$\overline T$ está definida sobre $K'$ significa que existe una parte cerrada
$T'$ de $X_{(K')}$ tal que $\overline T=g^{-1}(T')$~; como
$g^{-1}(X_{(K')}\setmin T')=X_{(K)}\setmin g^{-1}(T')$ para toda parte $T'$ de
$X_{(K')}$, \emph{a)}, \emph{b)} y \emph{b')} son equivalentes. Si $\overline Z$, cuyo espacio
subyacente es $\overline T$, está definido sobre $K'$, se tiene
$\overline Z=Z'\otimes_{K'}K$, donde $Z'$ es un subpreesquema cerrado de
$X_{(K')}$, y si $T'$ es el subespacio subyacente a $Z'$, se sabe que
$\overline T=g^{-1}(T')$ (\textbf I, 4.4.1), por lo que $\overline T$ está definida
sobre $K'$. Recíprocamente, para ver que \emph{a)} implica \emph{c)}, basta considerar un
subpreesquema cerrado $Z'$ de $X_{(K')}$ cuyo espacio subyacente sea $T'$
(\textbf I, 5.2.1) y tomar $\overline Z=Z'\otimes_{K'}K$.
\end{proof}

\begin{corollary}[4.9.2]
\label{IV.2.4.9.2-fr}
Con las notaciones de \textup{(4.9.1)}, supongamos que $K'$ sea un cuerpo de definición
de $\overline T$~; toda extensión $K''\subset K'$ de $k$ tal que $K'$ sea una
extensión radicial de $K''$ es también un cuerpo de definición de
$\overline T$.
\end{corollary}

\begin{proof}
Basta observar que la proyección canónica
$g:X_{(K')}\to X_{(K'')}$ es un morfismo entero, sobreyectivo y radicial, y
por tanto un homeomorfismo (2.4.5).
\end{proof}

\begin{remark}[4.9.3]
\label{IV.2.4.9.3-fr}
De (4.9.2) se deduce que una parte cerrada $\overline T$ de $X_{(K)}$ no
admite necesariamente un \emph{menor cuerpo de definición}. Por
ejemplo, si $K=k(t)$, donde $t$ es una indeterminada y $k$ un cuerpo perfecto
de característica $p>0$, una parte cerrada $\overline F$ de $X_{(K)}$ solo
puede tener un menor cuerpo de definición si ya está definida sobre
$k$: en efecto, para todo subcuerpo $K'$ de $K$ que contiene a $k$ y es
$\ne k$, se tiene $K'^p\ne K'$, $K'^p$ contiene a $k$ y $K'$

\oldpage[IV-2]{85}

es una extensión radicial de $K'^p$, por lo que, si $K'$ es un cuerpo de
definición de $\overline F$, lo mismo ocurre con $K'^p$. Sin embargo:
\end{remark}
